\chapter{Tinjauan Pustaka dan Dasar Teori}

\section{Tinjauan Pustaka}

Bagian ini membahas penelitian terdahulu yang relevan dengan implementasi I2S pada FPGA, integrasi peripheral kustom pada SoC berbasis prosesor soft-core, serta penggunaan AXI4-Lite sebagai antarmuka register.

\subsection{Penelitian tentang implementasi I2S pada FPGA}

Zhang dan Liu (2019) mengembangkan IP core I2S transmitter dan receiver untuk aplikasi audio processing pada FPGA Xilinx Zynq-7000. Penelitian tersebut menunjukkan bahwa antarmuka I2S dapat diimplementasikan secara andal pada perangkat logika terprogram dengan latensi rendah dan dukungan \textit{sample rate} hingga 96 kHz. Keterbatasannya adalah ketergantungan pada subsistem prosesor Zynq dan DMA, sehingga tidak secara langsung merepresentasikan peripheral kustom sederhana berbasis AXI4-Lite.

Anderson dan Lee (2021) merancang antarmuka I2S untuk \textit{digital microphone array} pada FPGA dengan fokus pada mode receiver dan perluasan kanal. Kontribusi penelitian ini terletak pada efisiensi berbagi resource pada beberapa kanal audio. Namun, ruang lingkupnya berbeda dengan penelitian ini karena tidak membahas transmitter audio menuju DAC eksternal.

Tanaka et al. (2020) mengimplementasikan subsistem audio digital pada FPGA dengan dukungan \textit{sample rate conversion}. Penelitian ini menegaskan pentingnya akurasi clock audio, buffering, dan sinkronisasi jalur data serial. Meskipun cakupannya lebih kompleks, hasil tersebut menunjukkan bahwa kualitas implementasi I2S sangat dipengaruhi oleh desain clock dan penanganan data antarblok.

\subsection{Penelitian tentang prosesor soft-core dan integrasi SoC}

Patterson dan Waterman (2017) memperkenalkan RISC-V sebagai ISA terbuka yang modular dan extensible, yang kemudian mendorong banyak implementasi prosesor soft-core pada FPGA. MicroBlaze V sebagai prosesor berbasis RISC-V dalam ekosistem AMD/Xilinx menjadi relevan karena menyediakan integrasi yang baik dengan Vivado dan Vitis.

Rodriguez et al. (2022) membandingkan performa beberapa konfigurasi prosesor soft-core pada FPGA dan menunjukkan bahwa arsitektur berbasis RISC-V sesuai untuk sistem embedded dengan peripheral memory-mapped. Hal ini mendukung pemilihan MicroBlaze V sebagai prosesor pengendali pada penelitian ini.

Sharma dan Kumar (2021) membahas pengembangan peripheral kustom menggunakan AXI4-Lite untuk sistem berbasis FPGA. Kontribusi utama penelitian tersebut adalah guideline praktis mengenai desain register map, integrasi dengan IP integrator, dan strategi pengembangan driver perangkat lunak sederhana.

\subsection{Penelitian tentang AXI4-Lite dan buffering data}

ARM (2011) melalui spesifikasi AMBA AXI4 menjadikan AXI4-Lite sebagai standar umum untuk peripheral dengan kebutuhan bandwidth rendah hingga menengah dan akses register sederhana. Sifatnya yang memory-mapped membuat peripheral mudah diakses dari perangkat lunak.

Chen et al. (2020) menganalisis overhead interkoneksi AXI pada SoC multi-peripheral dan menunjukkan bahwa AXI4-Lite cukup efisien untuk konfigurasi register, status, dan transfer data berukuran kecil. Untuk aliran data kontinu seperti audio, penelitian-penelitian tersebut juga mengindikasikan perlunya mekanisme buffering lokal agar perangkat lunak tidak harus melayani peripheral pada setiap sampel.

\subsection{Gap penelitian dan kontribusi}

Berdasarkan tinjauan di atas, dapat diidentifikasi beberapa celah penelitian:
\begin{enumerate}
	\item Banyak penelitian I2S pada FPGA berfokus pada subsistem audio yang lebih besar, DMA, atau platform prosesor tertentu, sehingga contoh peripheral I2S TX yang ringkas dan terintegrasi melalui AXI4-Lite masih relatif terbatas.
	\item Dokumentasi lengkap mengenai pengembangan \textbf{I2S TX custom IP} untuk \textbf{MicroBlaze V} masih sedikit, khususnya yang mencakup RTL, register map, integrasi Vivado, aplikasi Vitis, dan validasi hardware.
	\item Studi yang menekankan penggunaan \textbf{FIFO TX} dan \textbf{interupsi watermark} untuk menjaga kontinuitas audio pada sistem bare-metal sederhana juga belum banyak disajikan secara terstruktur.
\end{enumerate}

Kontribusi penelitian ini adalah:
\begin{enumerate}
	\item Merancang dan mengimplementasikan \textbf{IP core I2S transmitter} berbasis Verilog dengan antarmuka \textbf{AXI4-Lite} yang kompatibel dengan \textbf{MicroBlaze V}.
	\item Menyediakan alur lengkap dari desain RTL, integrasi SoC, pengembangan perangkat lunak uji, hingga validasi pada hardware Basys3 dengan PCM5102A.
	\item Menunjukkan penggunaan \textbf{FIFO TX} dan \textbf{watermark interrupt} sebagai strategi praktis untuk mengurangi underrun pada aliran audio.
	\item Menyajikan hasil implementasi dan pengujian yang dapat menjadi referensi untuk penelitian dan proyek audio digital berbasis FPGA.
\end{enumerate}

\section{Dasar Teori}

\subsection{Protokol komunikasi I2S}

I2S (Inter-IC Sound) adalah protokol serial sinkron yang dirancang untuk mengirimkan data audio digital antarperangkat. I2S digunakan secara luas untuk menghubungkan prosesor, FPGA, codec, DAC, dan ADC audio.

\subsubsection{Sinyal I2S}

I2S menggunakan tiga sinyal utama:
\begin{itemize}
	\item \textbf{Serial Data (SD)}: membawa data audio secara serial.
	\item \textbf{Word Select (WS/LRCLK)}: menandai kanal kiri dan kanan; frekuensinya sama dengan \textit{sample rate}.
	\item \textbf{Bit Clock (BCLK)}: clock serial yang menentukan perpindahan bit data.
\end{itemize}

\subsubsection{Format data Philips I2S}

Pada format Philips I2S:
\begin{itemize}
	\item data dikirim \textbf{MSB first};
	\item transisi \texttt{WS} terjadi satu periode clock sebelum MSB kanal berikutnya;
	\item data kiri dan kanan dikirim bergantian secara serial;
	\item implementasi praktis sering menggunakan slot 16, 24, atau 32 bit per kanal.
\end{itemize}

Pada penelitian ini, data audio dikirim sebagai \textbf{stereo 24 bit per saluran dalam slot 32 bit} sehingga satu frame stereo memerlukan 64 pulsa BCLK.

\subsubsection{Relasi \textit{sample rate}, BCLK, dan format slot}

Frekuensi sampling \(F_s\) menentukan jumlah frame audio stereo per detik. Jika satu frame menggunakan \(N\) bit clock, maka secara konseptual berlaku:
\begin{equation}
	F_s = \frac{f_{\mathrm{BCLK}}}{N}
\end{equation}

Untuk format stereo 32-bit per kanal, \(N = 64\). Karena itu, pengaturan BCLK dan pembagi clock internal sangat menentukan nilai \(F_s\) aktual.

\subsection{Audio digital, bit depth, dan DAC PCM5102A}

\textit{Sample rate} yang umum digunakan dalam sistem audio meliputi 44,1 kHz, 48 kHz, dan 96 kHz. Sementara itu, \textit{bit depth} menentukan resolusi amplitudo audio, misalnya 16 bit untuk audio konsumen dan 24 bit untuk audio berkualitas lebih tinggi.

PCM5102A adalah DAC audio stereo yang menerima data digital melalui antarmuka audio serial seperti I2S. Dalam konteks penelitian ini, PCM5102A berperan sebagai perangkat penerima data I2S dari FPGA. Keberhasilan sistem sangat bergantung pada ketepatan timing BCLK, WS, DATA, dan kesesuaian koneksi pin antara FPGA dan modul DAC.

\subsection{FIFO audio dan interupsi watermark}

Pada sistem yang dikendalikan prosesor, aliran sampel audio tidak selalu dapat dikirim tepat satu per satu pada waktu yang sama dengan kebutuhan serializer. Karena itu digunakan \textbf{FIFO TX} sebagai buffer antara prosesor dan logika serial I2S.

Jika level FIFO turun di bawah batas tertentu, peripheral dapat menghasilkan \textbf{interupsi watermark} agar prosesor segera mengisi ulang data. Pendekatan ini menurunkan frekuensi interupsi dibandingkan model satu interupsi per sampel dan membantu mencegah \textit{underrun}.

\subsection{Clocking Wizard dan frekuensi aktual}

Pada FPGA Xilinx, Clocking Wizard digunakan untuk menghasilkan clock internal yang mendekati frekuensi target. Akan tetapi, nilai yang dihasilkan sering berupa \textit{actual frequency} yang sedikit berbeda dari nilai nominal. Selisih kecil ini berpengaruh langsung pada frekuensi sampling terukur dan perlu dicatat saat evaluasi hasil.

\subsection{MicroBlaze V dan memory-mapped I/O}

MicroBlaze V adalah prosesor soft-core berbasis RISC-V yang terintegrasi dalam ekosistem Vivado dan Vitis. Peripheral kustom diakses melalui model \textbf{memory-mapped I/O}, yaitu register peripheral dipetakan ke alamat tertentu dalam ruang alamat prosesor.

Dengan pendekatan ini, perangkat lunak dapat:
\begin{itemize}
	\item menulis register kontrol untuk mengaktifkan peripheral;
	\item mengisi register data atau FIFO;
	\item membaca register status untuk memonitor kondisi sistem.
\end{itemize}

\subsection{AXI4-Lite}

AXI4-Lite adalah subset sederhana dari protokol AXI4 yang ditujukan untuk peripheral berbasis register. AXI4-Lite memiliki kanal alamat, data, dan respons untuk operasi baca/tulis tanpa dukungan burst transfer. Sifatnya yang sederhana menjadikannya cocok untuk peripheral I2S TX pada penelitian ini.

\subsection{Verilog HDL dan perancangan RTL}

Verilog digunakan untuk mendeskripsikan logika perangkat keras pada level RTL (\textit{Register Transfer Level}). Pada penelitian ini, Verilog digunakan untuk membangun:
\begin{itemize}
	\item register interface AXI4-Lite;
	\item logika pembagi clock;
	\item FIFO TX;
	\item serializer I2S;
	\item logika status dan interupsi.
\end{itemize}

\subsection{Analisis pemilihan metode}

Metode yang dipilih pada penelitian ini didasarkan pada beberapa pertimbangan:
\begin{enumerate}
	\item \textbf{MicroBlaze V} dipilih karena terintegrasi dengan baik pada ekosistem Vivado dan Vitis serta sesuai untuk sistem embedded bare-metal.
	\item \textbf{AXI4-Lite} dipilih karena sesuai untuk peripheral berbasis register dengan konfigurasi sederhana.
	\item \textbf{Verilog HDL} dipilih karena umum digunakan untuk desain RTL pada FPGA dan memiliki dukungan sintesis yang baik pada Vivado.
	\item Pendekatan \textbf{hardware-software co-design} dipilih agar peripheral dapat divalidasi tidak hanya dari sisi RTL, tetapi juga dari sisi penggunaan nyata oleh perangkat lunak.
\end{enumerate}
